\documentclass[11pt]{article}
\usepackage{amsmath, amssymb, amsthm, mathtools}
\usepackage[margin=1in]{geometry}

% Shortcuts
\newcommand{\N}{\mathbb{N}}
\newcommand{\Z}{\mathbb{Z}}
\newcommand{\Q}{\mathbb{Q}}
\newcommand{\R}{\mathbb{R}}
\newcommand{\set}[1]{\{#1\}}

% Theorem-like environments
\newtheorem{theorem}{Theorem}
\newtheorem{lemma}{Lemma}
\newtheorem{claim}{Claim}
\theoremstyle{definition}
\newtheorem{definition}{Definition}

\title{Problem Set 1}
\author{Michael Duren}
\date{\today}

\begin{document}
\maketitle

\section*{Problem 1}

\textbf{Claim.} Replace this with the statement you are proving.

\begin{proof}
  Write your proof here. Useful starting moves:
  \begin{itemize}
    \item \emph{Direct:} Assume the hypothesis, derive the conclusion.
    \item \emph{Contrapositive:} Assume $\neg Q$, derive $\neg P$.
    \item \emph{Contradiction:} Assume $P \land \neg Q$, derive a contradiction.
    \item \emph{Induction:} Base case $n = 0$, then $P(k) \implies P(k+1)$.
  \end{itemize}

  Example display equation:
  \[
    \sum_{i=1}^{n} i = \frac{n(n+1)}{2}.
  \]

  Conclude with what was shown. \qedhere
\end{proof}

\end{document}
